% !TeX root = ../main.tex
\begin{figure}[htbp]
\centering
\footnotesize
\textbf{Type grammar}\leanref{types}
\[
 \tau ::= \tyU\mid\tyB\mid\tyR m\mid\tau\to\tau
       \mid\tau\times\tau\mid\tau+\tau\mid\tyL\tau
 \qquad m ::= \G\mid\E
\]
\par\vspace{8pt}
\textbf{Real typing}\leanref{typing}
\begin{mathpar}
 \inferrule*[right=Real]{r\in\Real}{\typing r{\tyR m}}
 \and
 \inferrule*[right=Reject]{ }{\typing{\kw{reject}}\tau}
 \and
 \inferrule*[right=Uniform]{\typing a{\tyR m}\\\typing b{\tyR m}\\\alpha=\M\lor\alpha=m}
   {\typing{\draw{uniform}\alpha{a,b}}{\tyR m}}
 \\
 \inferrule*[right=Cmp]{\typing{e_1}{\tyR\G}\\\typing{e_2}{\tyR\G}}
   {\typing{e_1<e_2}{\tyB}}
 \and
 \inferrule*[right=Gaussian]{\typing u{\tyR m}\\\typing v{\tyR\G}\\\alpha=\M\lor\alpha=m}
   {\typing{\draw{gaussian}\alpha{u,v}}{\tyR m}}
 \\
 \inferrule*[right=Add]{\typing{e_1}{\tyR m}\\\typing{e_2}{\tyR m}}
   {\typing{e_1+e_2}{\tyR m}}
 \and
 \inferrule*[right=Mul]{\typing{e_1}{\tyR\G}\\\typing{e_2}{\tyR m}}
   {\typing{e_1\cdot e_2}{\tyR m}}
\end{mathpar}
\par\vspace{8pt}
\textbf{Subtyping}\leanref{subtyping}
\begin{mathpar}
 \inferrule*[right=Refl]{ }{\tau<:\tau}
 \and
 \inferrule*[right=Trans]{\tau<:\sigma\\\sigma<:\rho}{\tau<:\rho}
 \and
 \inferrule*[right=Real-Sub]{ }{\tyR\G<:\tyR\E}
 \and
 \inferrule*[right=Arrow]{\sigma_1<:\tau_1\\\tau_2<:\sigma_2}
   {\tau_1\to\tau_2<:\sigma_1\to\sigma_2}
\end{mathpar}
\caption{Type grammar, real typing, and subtyping rules. We show selected rules for arithmetic operations, distributions, and subtyping; additional rules for these constructs appear in~\Cref{sec:appendix-typing}. For distribution modes $\alpha\in\{\G,\E\}$, the type annotation $m$ must match $\alpha$; for $\alpha=\M$, either type annotation is permitted. We omit the standard rules for variables, unit and Boolean literals, functions,
recursion, application, let bindings, conditionals, pairs, projections, sum
injections and case analysis, and list constructors and case analysis.}
\label{fig:typing}
\Description{Type grammar, real typing rules, and subtyping rules. Subtyping permits G reals where E reals are expected. Multiplication requires a G left operand, division a G denominator, and comparison two G operands. Distribution rules share the same premises for sampling and mean evaluation.}
\end{figure}
